\documentclass[11pt]{article}
\usepackage[a4paper,margin=2.5cm]{geometry}
\usepackage{parskip}
\usepackage[T1]{fontenc}
\usepackage{lmodern}
\usepackage{microtype}
\usepackage{enumitem}

\setlist[itemize]{leftmargin=*,topsep=2pt,itemsep=2pt}

\begin{document}

{\LARGE\bfseries COMP3242/COMP6242 SCRIPT}

\vspace{4pt}

\noindent\textbf{Slide 1 \quad Introduction \quad 20 seconds}

Hi I'm Manindra and I'm here with Arjun, Razeen and Will and we will be presenting
Fine-Tuning of BioCLIP~2.5 with Taxonomic Heads for Multi-Species Plant
Identification.

\vspace{6pt}

\noindent\textbf{Slide 2 \quad Introduction -- Task + Motivation \quad 40 seconds}

Monitoring plant biodiversity at scale is essential for understanding ecosystem
health, tracking climate change, and informing conservation policy. Vegetation
surveys today are slow and expensive: identifying plant species in field
quadrats requires expert botanists on-site, and the supply of trained
taxonomists doesn't scale to the area that needs monitoring. We tackle this by
building a deep learning model that takes a single overhead quadrat photograph
and predicts every plant species present, drawn from a catalogue of 7,800
candidates. We're evaluated on 2,105 hidden test quadrats using macro-averaged
F1, averaged first within each transect and then across transects, so that
densely sampled sites don't dominate the final score.

\vspace{6pt}

\noindent\textbf{Slide 3 \quad Introduction -- Challenges \quad 40 seconds}

Two challenges make this much harder than standard classification. The first
is domain shift, illustrated in figure 1: training data is clean Pl@ntNet
close-ups of one plant in isolation, while test images are dense in-situ
quadrats where many species overlap and occlude one another. A model that
excels on isolated specimens has to generalise to crowded scenes it has never
been shown during training. The second is a severe long-tail: about twelve
species have more than a thousand training images, but 145 carry only a single
image, so standard cross-entropy training learns to over-predict common
species, adding false positives on quadrats where they aren't present and
missing the rare species that are. Both pull the per-quadrat F1 down, and
therefore the macro-averaged score.

\vspace{6pt}

\noindent\textbf{Slide 4 \quad Method Final Model \quad 60 seconds}

Our system has two stages. In the building stage, we take BioCLIP~2.5, a
vision model pretrained on 10 million organism images, and partially fine-tune
it on 2.65 million single-plant images. We freeze the lower 28 transformer
blocks to preserve its biological knowledge and only update the last 4. Three
auxiliary heads at the family, genus, and species levels regularise training,
though only the species head is used at inference. In the inference stage,
each quadrat is cut into a $4{\times}4$ grid of tiles. Each tile passes
through the trained model, probabilities are averaged, and an adaptive
threshold selects the final species set.

\vspace{6pt}

\noindent\textbf{Slide 5 \quad Ablations \quad 60 seconds}

Our most significant ablation was testing the number of unfrozen transformer
blocks in BioCLIP~2.5 to find the optimal number. As shown in the graph,
public Macro F1 peaks at eight unfrozen blocks at 0.391, then declines while
held-out validation accuracy kept rising. This finding showed that validation
accuracy on single-plant images does not predict quadrat performance. Full
fine-tuning collapsed to 0.301 as the model forgot its biological knowledge.
On our larger final dataset the optimal number of blocks was 4, suggesting
more data reduces the useful unfreeze depth. Beyond unfreezing, we also
ablated tile aggregation, logit adjustment, threshold selection, and resolution
ensembling on the final model.

\vspace{6pt}

\noindent\textbf{Slide 7 \quad Results \quad 40 seconds}

Other positive results included a softmax-mean tile aggregation, dual-resolution
ensembling at 224 and 336 pixels, logit adjustment to counteract long-tail
bias, and an adaptive threshold that always predicts between 2 and 10 species
per quadrat. Three additional pivots, a triple-backbone retrain, genus
co-occurrence reranking, and a seasonal phenology prior, all failed to
improve. Validation accuracy correlated weakly with quadrat leaderboard
performance at $r{\approx}0.4$, showing that the domain shift is the primary
bottleneck.

\vspace{6pt}

\noindent\textbf{Slide 8 \quad Further Research and Conclusion \quad 30 seconds}

Our best result was 0.418 public Macro F1 from a partially fine-tuned
BioCLIP~2.5 with tiled inference and dual-resolution ensembling, good for
7th place on the private leaderboard. Future gains may require training on
synthetic multi-species scenes that better match the test distribution.

\vspace{10pt}

\noindent\textbf{Marking Scheme}
\begin{itemize}
  \item Format and Communication (10 marks)
  \item Problem and Motivation (8 marks) (80 seconds $\sim$ 130 words)
  \item Method and Technical Content (12 marks) (120 seconds $\sim$ 200 words)
  \item Results and Conclusion (10 marks) (100 seconds $\sim$ 160 words)
\end{itemize}

\end{document}
